\section{Registro de plazos, captura histórica y consulta}\label{sec:matriz-plazos-persistentes}

La \cref{tab:matriz-plazos-persistentes} relaciona la implementación persistente con sus comprobaciones. La presencia de un caso en el repositorio identifica la propiedad exigida; su ejecución y resultado pertenecen al informe técnico. Los casos sintéticos no certifican reglas jurídicas, y la comprobación del transporte no sustituye la aceptación de una interfaz de usuario.

\begin{table}[H]
  \centering
  \caption{Fronteras de comprobación del seguimiento persistente de plazos.}
  \label{tab:matriz-plazos-persistentes}
  \small
  \begin{tabularx}{\textwidth}{p{3.2cm} X}
    \toprule
    \textbf{Frontera} & \textbf{Propiedad comprobada} \\
    \midrule
    Entrada \texttt{DEVI1} & Selección, condiciones, cantidad opcional y precisión reversibles; rechazo de representaciones inválidas. \\
    Resultado \texttt{DRES1} & Candidata, bloqueos y traza conservados sin recalcular el historial. \\
    Preparación y recibo & Vinculación de contenido revisado, contexto capturado y envío; ausencia de fechas fabricadas. \\
    Servicio y consultas & Permisos, reautenticación, ámbito, revisión exacta y cursores coherentes. \\
    PostgreSQL & Estado y auditoría atómicos, responsable elegible, comparación de base y fuentes vigentes. \\
    Historia y atención & Cálculo inmutable ante atención o retiro, con dependencias históricas independientes de las cabezas actuales. \\
    Seguimiento V2 & Observaciones verificadas, políticas explícitas, ambas huellas del predecesor y continuidad de decisiones humanas. \\
    Consumidor local & Revisión, resultado y auditoría atómicos; motivos sin cambios ligados a evidencia histórica exacta. \\
    Intentos y recuperación & Fallos tipados, espera persistida, conciliación de respuestas y conservación de todos los intentos. \\
    Catálogo y respaldo & Restricciones y proyecciones verificadas; restauración sujeta a contraste de todas las revisiones. \\
    HTTP & Comandos acotados y proyección estructurada de entradas, resultados, recibos e historia. \\
    \bottomrule
  \end{tabularx}
\end{table}

\begingroup
\raggedright
Los archivos se agrupan por frontera para mantener trazabilidad sin confundir pruebas auxiliares con campañas independientes. En el directorio \rutarepo{crates/application/tests/}, los nombres siguientes comparten el prefijo \texttt{deadline\_} y la extensión \texttt{.rs}:
\begin{itemize}
  \item Entrada: \texttt{evaluation\_input\_encoding}. Resultado y traza: los archivos cuyo sufijo comienza con \texttt{evaluation\_record}.
  \item Preparación y recibos: \texttt{preparation}, \texttt{receipts}, \texttt{receipt\_vectors} y \texttt{state\_bytes}.
  \item Autorización, contexto administrativo y consultas: \texttt{service}, \texttt{service\_administration} y \texttt{queries}.
\end{itemize}

Las pruebas de infraestructura residen en \rutarepo{crates/infrastructure/tests/}. También comparten el prefijo \texttt{deadline\_} y la extensión \texttt{.rs}. Sus sufijos identifican los grupos siguientes:
\begin{itemize}
  \item Los que comienzan con \texttt{backend}: operaciones, revalidación, concurrencia y familias de cálculo.
  \item \texttt{storage} e \texttt{input\_history}: reconstrucción exacta y corrupción.
  \item \texttt{input\_projection}, \texttt{schema}, \texttt{import} y \texttt{restore}: proyección SQL, arranque, importación y respaldo.
\end{itemize}
Las migraciones del grupo \texttt{0017} residen en \rutarepo{migrations/}.

En el directorio \rutarepo{crates/web/tests/}, los archivos con prefijo \texttt{deadline\_http} comprueban el contrato de transporte. Las pruebas de proyección del resultado se encuentran en \rutarepo{crates/web/src/deadlines/}, archivo \texttt{result\_tests.rs}. El informe \rutarepo{docs/verification-report.md} distingue pruebas focales, adaptadores reales y recorrido HTTP integrado. El documento \rutarepo{docs/academic-report-verification.md} registra por separado la compilación y revisión visual del manuscrito.
\par
\endgroup

La captura y la consulta de plazos en Qadra conservan el contrato V1. La composición del despachador y el consumidor en \texttt{serve}, el servicio humano y HTTP V2, las alertas y la agenda conjunta requieren sus recorridos de aceptación, junto con la validación de perfiles jurídicos aplicables. El alcance mantiene cálculos en días, meses y horas, sin reemplazarlos por captura manual de fechas.


\subsection{Selector e interfaz de plazos}

Las comprobaciones del selector mantienen el prefijo \texttt{deadline\_responsibles}. En \rutarepo{crates/application/tests/} se verifican servicio y respuestas; en \rutarepo{crates/web/tests/}, el contrato HTTP. Los archivos terminados en \texttt{backend.rs} y \texttt{concurrency.rs} del directorio de pruebas de infraestructura comprueban selección SQL y revocación concurrente.

En \rutarepo{web/tests/}, los archivos cuyo nombre comienza por \texttt{deadline-} comprueban el cliente, resultados, recibos, referencias y campos. El subdirectorio \texttt{browser/} agrega navegación, consultas, acceso, edición y presentación con transporte controlado. En \texttt{live/}, los archivos de plazos terminados en \texttt{workflow.spec.mjs}, \texttt{permissions.spec.mjs} y \texttt{concurrency.spec.mjs} ejercitan el servidor real. Se ejecutan mediante \rutarepo{scripts/web-demo.sh}, junto con los recorridos de regresión existentes.

El recorrido \rutarepo{scripts/api-demo.sh} incorpora la consulta de responsables y comprueba que la proyección no exponga credenciales. Los ensayos de navegador conservan resultados diarios, mensuales y horarios, relaciones históricas exactas, atención y retiro, cierre y revocación, conflictos y conciliación sin reenvío. La matriz separa el transporte controlado de las campañas reales; sus conteos y resultados observados permanecen en el informe de verificación.

\subsection{Seguimiento y consumidor durable}

\begingroup
\raggedright
Los valores y puertos están en \rutarepo{crates/application/src/}. Los módulos \texttt{deadline\_reevaluation}, \texttt{deadline\_observations} y \texttt{deadline\_technical} separan recibos, evidencia observada y preparación; \texttt{deadline\_worker} define resultados e intentos. En \rutarepo{crates/infrastructure/src/}, los módulos \texttt{deadline\_worker\_postgres} y \texttt{deadline\_worker\_schema} implementan consumo, lectura histórica y verificación de apertura. Las migraciones \texttt{0018} conservan V1/V2; las del grupo \texttt{0020} incorporan resultados, intentos y las restricciones de la escritura técnica.

En \rutarepo{crates/infrastructure/tests/}, los archivos comparten el prefijo \texttt{deadline\_worker\_} y la extensión \texttt{.rs}. Los sufijos \texttt{backend}, \texttt{guards} y \texttt{retries} identifican las once pruebas de la campaña focal inicial. \texttt{catalog} y \texttt{recovery} corresponden a las cinco pruebas focales de apertura, permisos, restauración real sin migración reparadora y crecimiento de reintentos. \texttt{profile\_failure} y \texttt{sql\_failures} corresponden a las tres pruebas focales de clasificación. Los archivos \texttt{families}, \texttt{history} e \texttt{inventory} amplían los escenarios del consumidor; sus casos se ejecutaron en la campaña integral, que incluye las 27 pruebas PostgreSQL del consumidor. El informe técnico conserva ese resultado sin volver a sumar las campañas focales.

El contrato está en \rutarepo{docs/deadline-worker.md}; el compromiso de observaciones y la continuidad, en \rutarepo{docs/deadline-tracking-receipts.md}; y el respaldo, en \rutarepo{docs/database-operations.md}. La decisión de arquitectura es \rutarepo{docs/adr/0037-durable-deadline-reevaluation.md}. El informe técnico conserva los resultados focales por separado y no los atribuye a HTTP o Qadra V2.
\par
\endgroup
